\documentclass[a4paper,11pt]{article}
%% Compiles both from this folder (pdflatex tp3.tex) and from the root of the project
%% (Overleaf: Menu > Main document = tp3_boosting/tp3.tex).
\makeatletter
\IfFileExists{tp3_boosting/tp3.tex}{\gdef\tpdir{tp3_boosting/}\gdef\tpcommon{common/}}{\gdef\tpdir{}\gdef\tpcommon{../common/}}
\edef\input@path{{\tpdir}{\tpcommon}}
\makeatother
\input{tpsetup}
\renewcommand{\polytitle}{Statistical Supervised Machine Learning -- TP 3}
\graphicspath{{figures/}{tp3_boosting/figures/}{../common/logos/}{common/logos/}}

\begin{document}
\sloppy
\sheethead{TP 3 -- Boosting: AdaBoost and gradient boosting}{%
Third of five practical sessions (1: protocol, 2: bagging, 3: boosting, 4 and 5: PAC-Bayes). Bagging (TP~2) averages
strong, unstable learners trained independently; boosting
builds a \emph{weighted} vote of weak learners trained \emph{sequentially}, each one focusing on the errors of the
previous ones. In this session we implement AdaBoost with decision stumps and gradient boosting with regression trees,
and we start to \emph{prove} things: short derivations explain why AdaBoost's weights have their form and why its
training error decreases exponentially fast. We then compare all the methods seen so far on real data, and we end by
looking at the vote of AdaBoost as a \emph{weighted} majority vote, with the margins and the quantities of Part~IV of the
lecture notes (Section~15). The
session lasts \textbf{three hours}; Part~\ref{p3:far} contains the difficult questions.}

\begin{sessionbox}[About this session -- TP 3: Boosting]
\textbf{Goal.} By the end of the session you should be able to implement AdaBoost and gradient boosting, explain the
role of the weights of the examples and of the voters, prove the exponential decrease of the training error of
AdaBoost, tune the learning rate and the number of rounds of gradient boosting, and compare ensemble methods with a
proper protocol.\par
\textbf{To reread before the session.} Section~10.2 (boosting, Algorithm AdaBoost and the empirical error bound of
AdaBoost) and Section~10.3 (gradient boosting) of the lecture notes; for Part~\ref{p3:far}, Section~15.6 (margins and
boosting).\par
\textbf{Timing (3 hours).}\quad
\begin{tabular}[t]{ll}Part~1 & 30 min\\
Part~2 & 45 min\\
Part~3 & 25 min\\
Part~4 & 45 min\\
Part~5 & 20 min\\
Part~6 \hard & 15 min, or at home\end{tabular}\par
\textbf{To hand in.} Your notebook with the tables of Parts~2, 4 and~5, the figures of Parts~2 and~4, and your
answers to Part~3 (on paper or in markdown).\par
\textbf{Allowed tools.} \code{numpy} for the stumps and the two boosting algorithms (regression trees of
\code{scikit-learn} may be used as base learners in Part~4); \code{scikit-learn} to check and to compare.
\end{sessionbox}

\medskip
\noindent\textbf{Data.} \code{wdbc}, \code{pima}, \code{australian} and \code{spambase} ($4597$ e-mails, $57$
features) loaded with \code{load(name)} of \code{common/tptools.py}, labels in $\{-1,+1\}$, and the \code{diabetes}
regression dataset of \code{scikit-learn}. The file \code{tp3\_starter.py} contains the skeletons.

\begin{tip}
Boosting runs hundreds of rounds, and each round learns a base classifier: anything that does not change from one
round to the next (here, the sort of each feature) must be computed once, before the loop. Store the scores
$f_t(x)=\sum_{s\leq t}\alpha_sh_s(x)$ of every round to plot the curves without recomputing anything.
\end{tip}

%% =====================================================================
\tppart{Decision stumps on weighted data}\label{p3:stumps}

A decision stump is a tree of depth one: $h_{j,t,s}(x)=s\cdot\sign(x_j-t)$, with a feature $j$, a threshold $t$ and a
sign $s\in\{-1,+1\}$. Given nonnegative weights $w_1,\dots,w_m$ summing to $1$, its weighted error is
$\varepsilon_w(h)=\sum_{i=1}^mw_i\ind{h(x_i)\neq y_i}$.

\begin{enumerate}
\item Let the examples be sorted by $x_j$, and let $L_k=\sum_{i\leq k}w_iy_i$ (sum over the $k$ first sorted examples)
      and $W=\sum_iw_iy_i$.\\
      Show that the stump with sign $s=+1$ cutting after the $k$-th example has weighted error
      $\frac12-\frac12(W-2L_k)$, and that the stump with $s=-1$ has error $1$ minus this quantity.
\item Write a class \code{Stumps(X)} which sorts every feature once in its constructor, and a method
      \code{fit(y, w)} which returns the best stump with no loop over the examples (use \code{np.cumsum}).
\item Check your method against a brute-force search (loops over $j$, $t$ and $s$) on $150$ examples of \code{pima}, with
      uniform and with random weights.\\
      How long do $100$ calls take on \code{spambase}?
\end{enumerate}

%% =====================================================================
\tppart{AdaBoost from scratch (Section 10.2)}\label{p3:ada}

\begin{enumerate}
\item Implement AdaBoost (Algorithm AdaBoost of Section~10.2) with your stumps: at round $t$, learn the stump $h_t$
      minimizing the weighted error $\varepsilon_t$, set $\alpha_t=\frac12\ln\frac{1-\varepsilon_t}{\varepsilon_t}$, update
      $w_i\leftarrow w_ie^{-\alpha_ty_ih_t(x_i)}/Z_t$, and return $H_T=\sign(f_T)$ with $f_T=\sum_{t\leq T}\alpha_th_t$.\\
      Store $\varepsilon_t$, $\alpha_t$, $Z_t$ and the scores $f_t$ on the training and test sets.
\item On \code{wdbc} and \code{spambase} ($70\%/30\%$), run $T=500$ rounds and plot, against $t$ (log scales), the
      training error, the test error, $\prod_{s\leq t}Z_s$ and $\exp\big(-2\sum_{s\leq t}(\frac12-\varepsilon_s)^2\big)$.
\item How do $\varepsilon_t$ and $\alpha_t$ evolve with $t$?\\
      Why?
\item Compare with \code{AdaBoostClassifier(DecisionTreeClassifier(max\_depth=1), n\_estimators=500)}.
\item What happens to the test error once the training error is $0$?\\
      Is this what you expected after TP~1 (overfitting) and TP~2 (number of trees in bagging)?
\end{enumerate}

%% =====================================================================
\tppart{Why AdaBoost works: three short proofs}\label{p3:why}

Let $w=w^{(t)}$ be the weights at round $t$, $h=h_t$, $\varepsilon=\varepsilon_t\in(0,\frac12)$ and, for $\alpha\in\R$,
$Z(\alpha)=\sum_iw_ie^{-\alpha y_ih(x_i)}$.

\begin{enumerate}
\item Show that $Z(\alpha)=(1-\varepsilon)e^{-\alpha}+\varepsilon e^{\alpha}$, that it is minimized by
      $\alpha_t=\frac12\ln\frac{1-\varepsilon}{\varepsilon}$, and that $Z(\alpha_t)=2\sqrt{\varepsilon(1-\varepsilon)}$.
\item Show that the weighted error of $h_t$ under the \emph{new} weights $w^{(t+1)}$ is exactly $\frac12$.\\
      Check it
      numerically.\\
      What does this imply for $h_{t+1}$?
\item Show that $w^{(T+1)}_i=\frac{e^{-y_if_T(x_i)}}{m\prod_{t\leq T}Z_t}$, and deduce, using
      $\ind{yf(x)\leq0}\leq e^{-yf(x)}$, that the training error satisfies $\RS(H_T)\leq\prod_{t\leq T}Z_t$.\\
      Then use
      $1-x\leq e^{-x}$ to obtain the bound of Section~10.2, $\RS(H_T)\leq\exp\big(-2\sum_t(\frac12-\varepsilon_t)^2\big)$.
\end{enumerate}

%% =====================================================================
\tppart{Gradient boosting for regression (Section 10.3)}\label{p3:gb}

Gradient boosting builds $F_M(x)=F_0+\nu\sum_{k=1}^Mg_k(x)$, where $F_0=\bar y$ and each $g_k$ is a small regression tree
fitted to the \emph{negative gradient} of the loss at the current predictions; $\nu\in(0,1]$ is the learning rate
(shrinkage).

\begin{enumerate}
\item For the squared loss $\ell(y,F)=\frac12(y-F)^2$, show that the negative gradient with respect to $F(x_i)$ is the
      residual $y_i-F_{k-1}(x_i)$.\\
      Write the algorithm.
\item Implement it with \code{DecisionTreeRegressor(max\_depth=2)} as base learner.\\
      On \code{diabetes}, split the data
      $60/20/20$ (training, validation, test).\\
      For $\nu\in\{1,0.3,0.1,0.03\}$ and up to $M=1000$ trees, plot the training
      and validation MSE against $M$.
\item For each $\nu$, select $M$ on the validation set (early stopping) and report the test MSE.\\
      Compare with
      \code{GradientBoostingRegressor} with the same parameters (\code{criterion="squared\_error"}).
\item What is the role of $\nu$?\\
      Why is boosting, unlike bagging, able to overfit when $M$ grows?
\end{enumerate}

%% =====================================================================
\tppart{Comparing the methods of the course}\label{p3:compare}

\begin{enumerate}
\item With $5$-fold stratified cross-validation, compare on \code{wdbc}, \code{pima}, \code{australian} and
      \code{spambase}: a logistic regression (with standardization in a pipeline), a tree of depth $5$, a random forest
      of $200$ trees, AdaBoost with $200$ stumps and gradient boosting with $200$ trees of depth $3$ (\code{scikit-learn},
      default values otherwise).\\
      Report the mean and the standard deviation of the error over the folds.
\item Is there a best method?\\
      Which differences seem significant, given the standard deviations?
\end{enumerate}

%% =====================================================================
\tppart{Going further \hard}\label{p3:far}

\begin{enumerate}
\item \textbf{AdaBoost minimizes the exponential loss.} Show that AdaBoost is a greedy, stage-wise minimization of
      $\frac1m\sum_ie^{-y_if(x_i)}$ over $f=\sum_t\alpha_th_t$: at round $t$, with $f_{t-1}$ fixed, the pair $(h,\alpha)$
      minimizing $\frac1m\sum_ie^{-y_i(f_{t-1}(x_i)+\alpha h(x_i))}$ is the stump of minimal weighted error and $\alpha_t$.
\item \textbf{Margins.} The normalised margin of a training example is $y_if_t(x_i)/\sum_{s\leq t}\alpha_s\in[-1,1]$.
      It is the margin $M_Q=1-2W_Q$ of Section~15.6, for the distribution $Q(h_s)=\alpha_s/\sum_{s'\leq t}\alpha_{s'}$ on
      the stumps.\\
      On \code{wdbc}, plot the cumulative distribution of the training margins for $t=10,50,500$, and relate it to the
      evolution of the test error after the training error has reached $0$.
\item \textbf{AdaBoost is a weighted majority vote.} With $Q(h_s)=\alpha_s/\sum_{s'\leq t}\alpha_{s'}$, $H_t=\BQ$.
      On \code{wdbc}, for $t\in\{10,100,500\}$, compute on the test set the Gibbs risk $\widehat R(\GQ)$, the joint error
      $\hat e(Q)$ and the risk of the vote, and the first-order and tandem quantities $2\widehat R(\GQ)$ and $4\hat e(Q)$
      (Propositions~15.2 and~15.6).\\
      Do the same for a random forest of $500$ trees (uniform $Q$).\\
      Which of the two votes do these bounds describe well? Why?
\item \textbf{Label noise.} On \code{spambase}, flip $20\%$ of the training labels (the test labels stay clean) and
      plot the test error of AdaBoost as a function of $t$ up to $1000$ rounds.\\
      Compare with the noiseless case and
      with a random forest.\\
      Explain using question~1.
\end{enumerate}

\end{document}
